% ======================================================================================
% Nom     : reports/latex/sections/04_pretraitement.tex
% Rôle    : Section décrivant le prétraitement des données.
% Auteur  : Maxime BRONNY
% Version : V1
% Cadre   : Réalisé dans le cadre du cours Techniques d'apprentissage artificiel
%           Master 1 Informatique Big Data
% Usage   : Fichier inclus dans main.tex lors de la compilation du rapport.
% ======================================================================================

\section{Prétraitement des données}
\label{sec:pretraitement}

Le prétraitement, qui s'appuie sur la bibliothèque pandas~\cite{mckinney2010pandas}
pour le chargement et la manipulation des données, est implémenté dans
\texttt{src/preprocessing.py} et \texttt{src/split\_data.py}. Un principe
structure l'ensemble :
\textbf{éviter toute fuite de données} (\emph{data leakage}). Tout paramètre
appris des données (médianes d'imputation, moyennes et écarts-types de
normalisation) est calculé \emph{uniquement} sur l'ensemble d'entraînement,
puis appliqué tel quel aux ensembles de validation et de test, qui jouent ainsi
pleinement leur rôle de données « jamais vues ».

\subsection{Nettoyage des valeurs aberrantes}

Les lignes physiquement impossibles identifiées lors de l'exploration sont
retirées : âge supérieur à 100 ans ou ancienneté d'emploi supérieure à 60 ans.
Seules \textbf{7 lignes sur 32\,581} sont concernées ($32\,574$ lignes
conservées). Nous ne touchons volontairement pas aux valeurs extrêmes mais
plausibles (revenus très élevés par exemple) : elles font partie de la vraie
distribution des données que les modèles doivent apprendre à gérer.

\subsection{Encodage des variables catégorielles}

Les quatre variables catégorielles sont converties en entiers par un
\emph{label encoding} à correspondance fixe, définie dans le code et non
apprise des données (tableau~\ref{tab:encodage}) ; l'encodage est donc
identique quel que soit le découpage.

\begin{table}[H]
  \centering
  \caption{Encodage des variables catégorielles.}
  \label{tab:encodage}
  \small
  \begin{tabular}{lp{8.5cm}}
    \toprule
    \textbf{Variable} & \textbf{Correspondance} \\
    \midrule
    \texttt{loan\_grade} & A$\to$0, B$\to$1, \ldots, G$\to$6 \\
    \texttt{cb\_person\_default\_on\_file} & N$\to$0, Y$\to$1 \\
    \texttt{person\_home\_ownership} & RENT$\to$0, OWN$\to$1, MORTGAGE$\to$2, OTHER$\to$3 \\
    \texttt{loan\_intent} & PERSONAL$\to$0, EDUCATION$\to$1, MEDICAL$\to$2,
      VENTURE$\to$3, HOMEIMPROVEMENT$\to$4, DEBTCONSOLIDATION$\to$5 \\
    \bottomrule
  \end{tabular}
\end{table}

Ce choix mérite justification. Pour \texttt{loan\_grade}, l'encodage entier est
naturel : la variable est \emph{ordinale} (A est un meilleur dossier que G) et
l'ordre numérique reflète un ordre réel. Pour
\texttt{cb\_person\_default\_on\_file}, binaire, l'encodage est sans ambiguïté.
Pour les deux variables \emph{nominales} (\texttt{person\_home\_ownership} et
\texttt{loan\_intent}), le label encoding introduit en revanche un ordre
artificiel entre modalités. Ce choix, hérité de la première version du projet
et conservé pour permettre la comparaison, reste acceptable pour l'arbre de
décision (qui peut isoler chaque modalité par divisions successives) et pour le
$k$-NN, mais constitue une limite pour la régression logistique ; nous y
revenons dans l'analyse critique (section~\ref{sec:analyse}).

\subsection{Découpage entraînement / validation / test}

Le jeu de données nettoyé est découpé en trois ensembles disjoints :
\textbf{70\,\% pour l'entraînement, 15\,\% pour la validation et 15\,\% pour le
test}, soit respectivement \textbf{22\,802, 4\,886 et 4\,886} observations.
Chaque ensemble a un rôle strictement distinct :

\begin{itemize}
  \item l'\textbf{entraînement} sert à apprendre les paramètres des modèles
    (poids de la régression, structure de l'arbre, instances mémorisées du
    $k$-NN) ;
  \item la \textbf{validation} sert \emph{uniquement} à choisir les
    hyperparamètres (seuil de décision, profondeur de l'arbre, nombre de
    voisins $k$) ;
  \item le \textbf{test} n'est utilisé qu'\emph{une seule fois}, à la toute
    fin, pour mesurer la performance de généralisation.
\end{itemize}

Cette séparation en trois est essentielle : si les hyperparamètres étaient
choisis directement sur le test, celui-ci ne mesurerait plus une capacité de
généralisation mais un ajustement indirect aux données d'évaluation, et les
performances rapportées seraient optimistes. La répartition $70/15/15$ offre un
bon compromis pour un jeu de cette taille : assez d'exemples d'entraînement
pour des modèles stables, et près de 5\,000 exemples dans chacun des deux
ensembles d'évaluation, ce qui donne des estimations fiables.

Le découpage est \textbf{stratifié} : il est effectué classe par classe, de
sorte que le taux de défaut ($21{,}8\,\%$) soit identique dans les trois
ensembles, comme le confirme la figure~\ref{fig:class-distribution}. Sans
stratification, un tirage aléatoire pourrait produire un ensemble de test non
représentatif. Le mélange utilise une graine aléatoire fixe ($42$), garantissant
la reproductibilité de toute l'expérience.

\begin{figure}[H]
  \centering
  \includegraphics[width=0.7\textwidth]{class_distribution.png}
  \caption{Taux de défaut par ensemble après découpage stratifié $70/15/15$ :
    les trois ensembles conservent la même proportion de défauts
    ($21{,}8\,\%$).}
  \label{fig:class-distribution}
\end{figure}

\subsection{Imputation des valeurs manquantes}

Les valeurs manquantes des variables numériques sont remplacées par la
\textbf{médiane} de la variable, calculée sur le seul ensemble d'entraînement.
La médiane est préférée à la moyenne car elle est robuste aux valeurs extrêmes,
nombreuses dans les variables de revenu et de montant.

\subsection{Normalisation}

Les variables sont enfin standardisées (score $z$) :
\[
  x' = \frac{x - \mu_{\text{train}}}{\sigma_{\text{train}}},
\]
où $\mu_{\text{train}}$ et $\sigma_{\text{train}}$ sont la moyenne et
l'écart-type calculés sur l'entraînement. Cette normalisation est indispensable
pour deux des trois modèles : pour la régression logistique, elle conditionne
la vitesse et la stabilité de la descente de gradient ; pour le $k$-NN, elle
évite qu'une variable à grande échelle comme le revenu (de l'ordre de
$10^4$ à $10^5$) n'écrase toutes les autres dans le calcul des distances.
L'arbre de décision y est insensible (il ne compare que des seuils par
variable), mais reçoit les mêmes données afin de conserver un pipeline unique.
